\documentclass[11pt]{article}
\usepackage[margin=1in]{geometry}
\usepackage{amsmath,amssymb,booktabs,tabularx}
\usepackage[T1]{fontenc}
\usepackage{hyperref}
\hypersetup{colorlinks=true,urlcolor=blue,citecolor=blue}
\newcommand{\pending}{\textit{Pending}}
\title{VLAQuantShift: A Protocol for Paired Quantization--Shift Interaction Audits}
\author{Research protocol draft\\Authorship and affiliations to be finalized}
\date{October 16, 2026}

\begin{document}
\maketitle
\begin{abstract}
This protocol specifies an extension of VLAQuantBench for evaluating whether
quantization changes a vision--language--action policy's sensitivity to
observation shifts. A paired factorial design compares released-precision and
quantized policies under clean and shifted observations. Its primary outcome is
the increase in the quantization penalty under a prespecified shift, with
episode identity preserved in task-clustered uncertainty estimates. Calibration,
precision selection, and evaluation use separate provenance manifests. A
secondary offline diagnostic decomposes supplied cross-policy/state action
arrays into reference-state discrepancy, state-response, and interaction terms.
The current artifact is a protocol and software candidate, not a completed
empirical study. New model-performance, real-robot, and deployment-speed results
are pending. The proposed contribution concerns reproducible interaction audits;
neither the factorial contrast nor the algebraic action identity is new theory.
\end{abstract}

\section{Research question and provenance}
Can a precision configuration preserve clean task success while increasing
degradation under observation changes? Does its ranking transfer across
calibration sources and execution horizons? These are hypotheses to test, not
findings of this draft.

The implementation develops VLAQuantBench~\cite{xu2026vlaquantbench}, starting
from upstream commit
\texttt{4a2cb7c1bf9160ae683fe598ec583d7c9b9c13c9}.
Its source, license, attribution, and historical results remain upstream work.
Model adapters, quantization routines, and evaluators inherited from that project
must be identified as such. The new candidate artifact adds a shift-evaluation
protocol and supporting analysis. Historical upstream measurements are not
relabeled as measurements from this study.

\section{Related work and overlap}
VLAQuantBench already studies closed-loop precision configurations, calibration
controls, and fixed-observation replay~\cite{xu2026vlaquantbench}.
QuantVLA calibrates attention and output scales~\cite{zhang2026quantvla};
QVLA measures action-sensitive channels for bit allocation~\cite{xu2026qvla}.
PIVOT-Q investigates quantized-policy generalization under environmental shifts
and selective correction of vulnerable states~\cite{ruan2026pivotq}.
DA-PTQ targets motion drift~\cite{xu2026daptq}, while DyQ-VLA adapts precision
temporally~\cite{zheng2026dyqvla}. ActQuant and HBVLA further establish
action-guided low-bit compression~\cite{akbari2026actquant,yan2026hbvla}.
Cross-layer dependencies in mixed precision are also established by
CLADO~\cite{deng2025clado}.

LIBERO-Plus supplies systematic robustness perturbations~\cite{fei2025liberoplus},
and VLABench supplies a broader manipulation setting~\cite{zhang2024vlabench}.
A public on-policy calibration repository reports paired experiments and a
negative calibration result~\cite{parastoopil2026calibrate}; it is cited as a
research software artifact, not as independently verified peer-reviewed evidence.
Consequently, closed-loop evaluation, action sensitivity, trajectory drift,
robustness, and on-policy calibration are not claimed as new ideas here.
Publication-level differentiation will require new, reproducible empirical
evidence beyond these overlaps.

\section{Paired factorial protocol}
Fix a checkpoint, quantization recipe, selected modules, execution convention,
and shift definition. For task $j$ and episode identity $i$, obtain binary
success $Y^{p,c}_{ji}$, where $p\in\{f,q\}$ denotes released precision or
quantization, and $c\in\{0,s\}$ denotes clean or shifted observations.
Each identity contributes all four conditions. Define
\begin{align}
e_{ji} &= (Y^{f,s}_{ji}-Y^{q,s}_{ji})
          -(Y^{f,0}_{ji}-Y^{q,0}_{ji}),\\
\widehat E &= \frac{1}{J}\sum_{j=1}^{J}
                 \frac{1}{n_j}\sum_{i=1}^{n_j}e_{ji}.
\end{align}
Positive $100\widehat E$ indicates additional quantization penalty under the
shift, in percentage points. All four success rates and both precision penalties
must accompany this statistic: shared FP/Q failure can yield zero interaction.
If shifted FP and Q success both reach zero, the interaction becomes the negative
clean precision penalty, even though neither shifted policy succeeds. Clean
ceilings similarly hide action-quality differences. Report paired outcome
patterns and task-level rates without removing floor/ceiling tasks post hoc.
This is a standard factorial contrast conditional on the specified intervention,
not a new estimator or an unrestricted claim of real-world causality.

\paragraph{Proposed primary test.}
Before final evaluation, freeze $\pi_{0.5}$ (\texttt{lerobot/pi05\_libero\_finetuned\_v044}) on
LIBERO-Spatial, end-to-end RTN W4, evaluation seed 0, and
\texttt{brightness:0.2} as the primary contrast, together with exact checkpoint
revision, module list, transform, task subset and code revisions. The transform
adds $0.2\times255$ to each uint8 RGB channel, clips to $[0,255]$ and casts to
uint8; it is not multiplicative brightness. Seeds 1 and 2 and OpenVLA-OFT are
secondary replications. A named model repository is not an immutable checkpoint
revision, and the current plan does not itself resolve one.
The primary metric is the equal-task estimate above, with a two-sided 95\%
task-cluster bootstrap interval. This is a preregistration proposal until the
manifest is committed; changes after viewing test outcomes are exploratory.

\paragraph{Observation intervention.}
The initial implementation applies deterministic observation transformations
through inherited runners. Specify camera coverage, dtype, pixel range, clipping,
ordering and random-seed derivation before preprocessing. Image corruption does
not reproduce a physical scene intervention. These transforms are not described
as an official LIBERO-Plus implementation. Scene randomization is a later,
separately validated integration.

\paragraph{Pair identity.}
Retain benchmark, suite, task ID, initial-state ID, seed, replicate and shift
metadata. Duplicate or incomplete four-cell records fail validation by default.
Record execution errors separately from unsuccessful completed episodes.
Matched initial states and RNG seeds do not alone guarantee matched per-step
noise: each policy adapter must disclose deterministic execution, explicit common
random numbers, or weaker seed-only matching. Action queues, caches and random
number consumption need adapter-specific audits.
Observation-corruption randomness uses a separate identity-derived seed. The
sample plans hold this seed fixed across policy seeds; they do not independently
replicate corruption realizations. Severity levels also need not share nested
noise or masks.

\section{Calibration separation and statistical analysis}
The final study requires immutable calibration, selection and evaluation
manifests. Enforce episode-disjointness for all roles and task-disjointness when
claiming task transfer. Current software checks calibration/evaluation episode
index overlap; selection manifests and task-disjoint validation are not yet
implemented. Quantizer parameters remain frozen across clean and shifted evaluation.
Shifted recalibration is a separate adaptation arm at a declared sample and
collection budget. Do not fit normalization or select precision from test data.

The primary interval resamples task clusters with replacement, carrying complete
paired episodes, for at least 5,000 seeded bootstrap replicates (override the
current CLI default of 2,000). Current software resamples tasks only, conditional
on recorded episodes, and reports one contrast per evaluation seed. It provides
neither a within-task resampling stage nor a pooled multi-seed interval.
Few clusters limit inference; provide all task-level effects. A collapsed interval
is not evidence of certainty. For calibrated methods, use at least three separate
calibration sets and report their results separately. Current configurations
couple calibration and evaluation RNG seeds and reuse the same calibration
initial states, so their three seed runs are not an independently crossed
calibration-by-evaluation design. Independent calibration sets require separate
plans with different episode ranges. These intervals do not account for new
within-task sampling, training uncertainty or calibration variation. Repeated
evaluation seeds sharing an initial condition are not independent initial states.

The cache key distinguishes calibration transforms and execution settings.
The extension records and compares calibration-cache content SHA256 values.
Checkpoint content digests still require separate retention and verification.
Current calibration collects unquantized-policy trajectories under one transform;
mixed calibration sources and quantized-policy collection remain extensions.

Secondary comparisons are descriptive unless a multiplicity procedure is
prespecified. An interval containing zero does not establish equivalence. Report
null effects, negative interactions, and calibration-seed variability.

\section{Offline cross-policy/state action diagnostic}
Let $s_f$ and $s_q$ represent the histories and states of two paired policy
branches at a common control index. Let
\[
a_{ff}=\pi_f(s_f),\quad a_{qf}=\pi_q(s_f),\quad
a_{fq}=\pi_f(s_q),\quad a_{qq}=\pi_q(s_q).
\]
Under specified policy memory and inference randomness, define
\begin{align}
D &= a_{qf}-a_{ff}, & S &= a_{fq}-a_{ff},\\
I &= a_{qq}-a_{fq}-a_{qf}+a_{ff}, & T &= a_{qq}-a_{ff}.
\end{align}
Then $T=D+S+I$ exactly. The terms describe reference-history quantization
discrepancy, the FP response to changed history, and their interaction. Norms do
not add; the corresponding energy identity is
\[
\|T\|^2=\|D\|^2+\|S\|^2+\|I\|^2
+2\langle D,S\rangle+2\langle D,I\rangle+2\langle S,I\rangle.
\]
Report closure error, component norms and cross terms. Fix action units and any
calibration-fitted normalization. Mask invalid timesteps after either trajectory
ends and disclose coverage rather than zero-padding. Aggregate within episodes
before across tasks. The current analyzer accepts finite equal-shaped $T\times D$
arrays; callers must prepare alignment, censoring and episode provenance, which
the analyzer cannot establish from the arrays alone.

The initial analyzer accepts offline arrays. Simulator snapshot/restore, paired
state branching and automatic crossed VLA evaluation are \emph{not implemented}
by this diagnostic. Future collection must audit recurrent state, action caches,
observation history, flow noise and action-chunk semantics. An algebraic identity
alone does not isolate a physical mechanism or establish that a discrepancy
caused task failure.

\section{Planned study and ablations}
Table~\ref{tab:plan} specifies intended coverage, not completed experiments.
Hold out tasks used for selection and a separate pilot; freeze eligible task
identities before evaluation. The initial target is at least 20 initial states
per held-out task and three evaluation seeds where feasible, with independent
calibration sets for calibrated controls. These are planning targets, not a power
guarantee. Reassess sample requirements with a separate pilot.

\begin{table}[ht]
\centering\small
\begin{tabularx}{\linewidth}{lX}
\toprule
Axis & Planned coverage \\
\midrule
Models & OpenVLA-OFT and $\pi_{0.5}$; pinned released checkpoints \\
Suites & LIBERO-Spatial and LIBERO-Long; VLABench expansion later \\
Precision & FP, RTN W4, W4A8, protected-output reference configuration \\
Shifts & Identity, additive brightness, seeded noise, occlusion, camera dropout; contrast later \\
Calibration & None for RTN; separate clean/shifted FP collections; mixed sources later \\
Method comparison & Reproduce available VLA PTQ methods before matched comparisons \\
Diagnostics & Paired outcomes first; supplied-array decomposition; branch collection later \\
\bottomrule
\end{tabularx}
\caption{Planned matrix. Observation shifts are distinct from simulator scene changes.}
\label{tab:plan}
\end{table}

Required controls include identity shifts, independent FP repeats, quantized
module/hook coverage, equal-budget calibration-source comparisons, protected
output modules, and executed action-count ablations. Predicted chunk length and
executed action count are separate variables. Precision selection uses selection
data only; an action-error selector is a baseline in light of QVLA and ActQuant.

\section{Results status and reporting contract}
No new VLA or physical-robot results are reported in this draft. Toy fixtures
validate software behavior only. Fill Table~\ref{tab:pending} from immutable
episode logs after the study; every cell must link to its manifest and records.

\begin{table}[ht]
\centering\small
\begin{tabular}{lccccc}
\toprule
Model / suite & FP clean & Q clean & FP shift & Q shift & Excess [95\% CI] \\
\midrule
OpenVLA-OFT / Spatial & \pending & \pending & \pending & \pending & \pending \\
$\pi_{0.5}$ / Spatial & \pending & \pending & \pending & \pending & \pending \\
OpenVLA-OFT / Long & \pending & \pending & \pending & \pending & \pending \\
$\pi_{0.5}$ / Long & \pending & \pending & \pending & \pending & \pending \\
\bottomrule
\end{tabular}
\caption{Planned reporting table. Pending entries are not zero or failed runs.}
\label{tab:pending}
\end{table}

Additional planned tables report per-task effects, calibration replicates,
precision-ranking transfer, missing records, and action-diagnostic coverage.
Fake quantization supports numerical accuracy evaluation; estimated bit budgets
do not establish packed memory, measured VRAM, kernel latency or robot speed.
Any deployment measurement requires a separate documented implementation and
device-specific experiment. GPU execution in the initial workspace is restricted
to CUDA devices 6 and 7.

\section{Limitations and completion criteria}
Observation corruptions cover only part of deployment variation. Checkpoints,
task sets, calibration sets and simulator conventions constrain generalization.
Stateful policies complicate crossed evaluation, and the initial release has no
automatic simulator-branch collector. The available methods and baseline
reproductions may limit the comparison. The study has strong nearby prior work,
especially PIVOT-Q and DA-PTQ, and novelty is not guaranteed.

An empirical manuscript requires validated adapters, a frozen preregistration,
complete paired model experiments, reproduced method controls, audited
uncertainty and new scientific evidence. Passing software tests or publishing
this repository does not satisfy those requirements.

\bibliographystyle{plain}
\bibliography{shift_references}
\end{document}
